\documentclass[10pt,a4paper]{amsart}
\usepackage[margin=19mm]{geometry}
\usepackage{amsmath,amssymb,mathtools}
\usepackage[scr=rsfs,cal=euler]{mathalfa}
\usepackage{xfrac}
\usepackage[unicode,hidelinks,bookmarks=false]{hyperref}
\newcommand{\ie}{i.e.\ }
\newcommand{\cf}{cf.\ }
\newcommand{\resp}{respectively\ }
\newcommand{\Subsection}{\subsection}
\providecommand{\nobreakdash}{\nobreak\hbox{-}}
\setlength{\emergencystretch}{2em}
\multlinegap0pt
\usepackage{bm}
\usepackage{bbm}
\usepackage{dsfont}
\usepackage{xspace}
\usepackage{cancel}
\usepackage{mathtools}
\usepackage{amsthm}
\newtheorem{theorem}{Theorem}
\newtheorem{proposition}[theorem]{Proposition}
\title[The mod 2 Seiberg-Witten invariants of spin structures and s]{The mod 2 Seiberg-Witten invariants of spin structures and spin families}
\author{David Baraglia}
\date{Source: arXiv:2303.06883v3 (CC BY 4.0)}
\begin{document}
\maketitle

\noindent\emph{Source and licence.} The mod 2 Seiberg-Witten invariants of spin structures and spin families. Version arXiv:2303.06883v3 (CC BY 4.0), distributed under the Creative Commons Attribution 4.0 International licence, as recorded in the arXiv licence metadata for this exact version and shown on the abstract page https://arxiv.org/abs/2303.06883v3. Full rights notice: licenses/arxiv-2303.06883-CC-BY-4.0.txt.\par\smallskip

\noindent\textit{Editorial scope.} Counted unit: the Proposition whose source label is \texttt{prop:q12} in the article (source0.tex lines 653--655, source label \texttt{prop:q12}), delivered together with the article's own complete original proof of it (source0.tex lines 656--673, one \texttt{proof} environment of 18 lines). No mathematics is written, repaired or re-derived: statement and proof are the article's own text line for line. Required context retained uncounted: none. Every definition, display and label of those lines is retained unchanged. Omitted from this package and not redistributed: 1--652, 674--1512. Nothing inside a retained excerpt is omitted or altered: every retained body is delivered exactly as the article's lines are written, so no omission receipt of any kind accompanies this package. Citations: the retained excerpts carry no citation command at all, so no bibliography entry of the article is transcribed and no reference is redistributed. Reference adaptation: none was needed and none was made; every reference inside the delivered unit resolves inside it, so no unresolved reference or citation is produced. Printed slips kept literal: no printed word, symbol, hypothesis or number of the counted statement or of its proof was altered. Layout and class: the article's own class is \texttt{amsart} with packages including amsfonts, amscd, amsthm, amsgen, amsmath, amssymb, xy, epsfig, color, youngtab; the wrapper is the lane's pinned \texttt{amsart} editorial wrapper at 10pt on A4, which restates the theorem-like environments the retained text needs and adds the article's own macro definitions that the retained text needs (none), with extra wrapper packages (bm, bbm, dsfont, xspace, cancel, mathtools). The delivered numbering is the unit's own. Assets: no retained span contains a figure environment, a graphics-inclusion command, a PDF-page inclusion, a table or a TikZ picture; no image file is copied into this package, the package declares no asset dependency and no image file is redistributed with it. Licence: version arXiv:2303.06883v3 of this article is distributed under the Creative Commons Attribution 4.0 International licence, as recorded in the arXiv licence metadata for this exact version and shown on the abstract page https://arxiv.org/abs/2303.06883v3. A full rights notice is delivered at licenses/arxiv-2303.06883-CC-BY-4.0.txt. Characters: every retained excerpt is pure ASCII, so the delivered engine, XeLaTeX with its default Latin Modern text font, reports no missing character.
\begin{proposition}\label{prop:q12}
We have $H^*_G(pt) \cong \mathbb{Z}_2[u,y,q_1]/(u^3)$. Furthermore we have $q_2 + q_1 = y^2 + yu^2$.
\end{proposition}

\begin{proof}
We have a short exact sequence $1 \to S^1_1 \to G \buildrel p_1 \over \longrightarrow Pin(2) \to 1$, where $S^1_1$ denotes the subgroup $S^1 \times \{1\} \subset S^1 \times S^1 \subset G$. The Lyndon--Hochschild--Serre spectral sequence for $H^*_G(pt)$ has $E_2^{*,*} = H^*_{Pin(2)}( H^*_{S^1_1}(pt) ) \cong H^*_{Pin(2)}[y']$, where $deg(y') = 2$ is the generator of $H^2_{S^1_1}$. The composition $S^1_1 \to G \buildrel p_1 \over \longrightarrow O(2)$ is easily seen to be the inclusion of the circle subgroup of $O(2)$. This implies that $p^*(y) \in H^2_G(pt)$ restricts to $y'$. This implies that the spectral sequence degenerates and $H^*_G(pt) \cong H^*_{Pin(2)}[y] \cong \mathbb{Z}_2[u,y,q_1]/(u^3)$.

It remains to prove the relation $q_1 + q_2 = y^2 + yu^2$. Since $H^4_{G}(pt)$ is spanned by $q_1,y^2,yu^2$, we have that $q_2$ is a linear combination of $q_1, y^2, yu^2$. Consider the subgroup $S^1 \times S^1 \subset G$ which has cohomology $\mathbb{Z}_2[x_1,x_2]$. The restriction map $H^*_{G}(pt) \to H^*_{S^1 \times S^1}$ sends $q_i$ to $x_i^2$, $y$ to $x_1+x_2$ and $u$ to zero. This shows that $q_2$ must be either $q_1 + y^2$ or $q_1 + y^2 + yu^2$. 

Next, we note that $Pin(2)$ acts freely on $S^3 \cong SU(2)$ via the inclusion $Pin(2) \to SU(2)$ and the left action of $SU(2)$ on itself. The quotient space is $\mathbb{RP}^2$, because the quotient of $S^3$ by the $S^1$-subgroup of $Pin(2)$ is $S^3/S^1 \cong S^2$ and the remaining action of $Pin(2)/S^1 \cong \mathbb{Z}_2$ acts on $S^2$ as the antipodal map.

We also have that $G$ acts freely on $S^3 \times S^3$ via the inclusion $G \to Pin(2) \times Pin(2)$ and the obvious product action of $Pin(2) \times Pin(2)$ on $S^3 \times S^3$. Let $M = (S^3 \times S^3)/G$ be the quotient. Clearly $M$ is the quotient of $S^2 \times S^2$ by the involution which acts as the antipodal map on both factors. Projecting to either factor of $S^2$ gives two fibrations
\[
S^2 \to M \buildrel \pi_i \over \longrightarrow \mathbb{RP}^2, \quad i = 1,2.
\]
Both fibrations admit a section $s_i : \mathbb{RP}^2 \to M$ which the image under the quotient map $(S^2 \times S^2) \to M$ of the diagonal copy of $S^2$. From this it follows easily that the Leray--Serre spectral sequences for both fibrations degenerate at $E_2$. Then $H^*_G( S^3 \times S^3) \cong H^*(M)$ is a $6$-dimensional space over $\mathbb{Z}_2$ with basis $1,u,u^2,c,cu,cu^2$, where $c \in H^2(M)$ restricts non-trivially to the fibres of $\pi_1 : M \to \mathbb{RP}^2$. The diagonal $S^2 \to S^2 \times S^2$ has normal bundle $TS^2$ and taking the quotient by the antipodal map on both factors, it follows that the normal bundle of the sections $s_1,s_2$ are both equal to $T\mathbb{RP}^2$. Since $w_2( T\mathbb{RP}^2) = u^2$, it follows that the mod $2$ self-intersection of $s_i$ is odd, but this can only happen if $c^2 \neq 0$, so $c^2 = cu^2$.

Let $\mathbb{H}_1$ be the standard representation of the $i$-th copy of $SU(2)$ in the product $SU(2) \times SU(2)$. By restriction, this defines a representation of $G$ and we have that $q_i = w_4( \mathbb{H}_i)$. By taking the pullback of $\mathbb{H}_i$ to $S^3 \times S^3$ and quotienting by $G$, we have that $\mathbb{H}_i$ descends to a rank $4$ vector bundle $\widetilde{\mathbb{H}}_i \to M$. In fact, $\widetilde{\mathbb{H}}_i$ is the pullback under $\pi_i : M \to \mathbb{RP}^2$ of a rank $4$ vector bundle on $\mathbb{RP}^2$, because the $G$-action on $\mathbb{H}_i$ factors through $p_i : G \to Pin(2)$. Now since $\mathbb{RP}^2$ is $2$-dimensional, we must have that $w_4(\widetilde{\mathbb{H}}_i) = 0$. So the pullback of $q_1,q_2$ to $H^4_{G}(S^3 \times S^3) \cong H^4(M)$ are both zero.

Let $R$ denote the standard $2$-dimensional real representation of $O(2)$. We have that $y = w_2(R)$. By taking the pullback of $R$ to $S^3 \times S^3$ and quotienting by $G$, we obtain a vector bundle $\widetilde{R} \to M$ on $M$. The restriction of $\widetilde{R}$ to any fibre of $\pi_2$ is isomorphic to $\mathcal{O}(1) \to S^2$. In particular, this means that $y = w_2(\widetilde{R})$ equals either $c$ or $c+u^2$. In either case, we then have $y^2 = cu^2 = yu^2$. So $y^2 + yu^2 = 0$, but $y^2 \neq 0$ in $H^*(M)$. Then since $q_1 + q_2 = 0$ in $H^*(M)$, we see that we must have $q_1 + q_2 = y^2 + yu^2$.

\end{proof}

\end{document}
